\documentclass[12pt]{article}
\usepackage[a4paper, margin=1in]{geometry}

\begin{document}

\begin{center}
    \Large \textbf{Kinetic Energy Worksheet}
\end{center}

\vspace{0.5cm}

\noindent \textbf{Formula:} \quad $E = \frac{mv^2}{2}$

\vspace{0.5cm}

\begin{enumerate}

\item A 2 kg object moves at 3 m/s. Calculate its kinetic energy.

\item An object has a mass of 5 kg and moves at 4 m/s. Calculate its kinetic energy.

\item An object has a kinetic energy of 100 J and a mass of 4 kg. Calculate its velocity.

\item An object has a kinetic energy of 72 J and moves at 6 m/s. Calculate its mass.

\item An object of mass 10 kg has a kinetic energy of 200 J. Calculate its velocity.

\item A car of mass 1000 kg is travelling at 20 m/s. Calculate its kinetic energy.

\item A football of mass 0.5 kg is kicked at 10 m/s. Calculate its kinetic energy.

\item A moving object has a kinetic energy of 150 J and a mass of 3 kg. Calculate its velocity.

\item A cyclist and bike have a combined mass of 80 kg and a kinetic energy of 640 J. Calculate their speed.

\item A ball moving at 8 m/s has a kinetic energy of 128 J. Calculate its mass.

\item A 500 g object is moving at 6 m/s. Calculate its kinetic energy.

\item A moving object has a kinetic energy of 2 kJ and a mass of 8 kg. Calculate its velocity.

\item A 250 g object has a kinetic energy of 50 J. Calculate its velocity.

\item A car has a kinetic energy of 5 kJ and is travelling at 10 m/s. Calculate its mass.

\item A 1200 kg car increases its speed from 10 m/s to 20 m/s. Calculate the increase in kinetic energy.

\end{enumerate}

\newpage

\begin{center}
    \Large \textbf{Answer Key}
\end{center}

\vspace{0.5cm}

\begin{enumerate}

\item $E = \frac{1}{2}(2)(3^2) = 9 \, J$

\item $E = \frac{1}{2}(5)(4^2) = 40 \, J$

\item $100 = \frac{1}{2}(4)v^2 \Rightarrow v = 7.07 \, m/s$

\item $72 = \frac{1}{2}m(6^2) \Rightarrow m = 4 \, kg$

\item $200 = \frac{1}{2}(10)v^2 \Rightarrow v = 6.32 \, m/s$

\item $E = \frac{1}{2}(1000)(20^2) = 200000 \, J$

\item $E = \frac{1}{2}(0.5)(10^2) = 25 \, J$

\item $150 = \frac{1}{2}(3)v^2 \Rightarrow v = 10 \, m/s$

\item $640 = \frac{1}{2}(80)v^2 \Rightarrow v = 4 \, m/s$

\item $128 = \frac{1}{2}m(8^2) \Rightarrow m = 4 \, kg$

\item $m = 0.5 \, kg \Rightarrow E = \frac{1}{2}(0.5)(6^2) = 9 \, J$

\item $E = 2000 \, J \Rightarrow 2000 = \frac{1}{2}(8)v^2 \Rightarrow v = 22.36 \, m/s$

\item $m = 0.25 \, kg \Rightarrow 50 = \frac{1}{2}(0.25)v^2 \Rightarrow v = 20 \, m/s$

\item $E = 5000 \, J \Rightarrow 5000 = \frac{1}{2}m(10^2) \Rightarrow m = 100 \, kg$

\item $\Delta E = \frac{1}{2}(1200)(20^2 - 10^2) = 180000 \, J$

\end{enumerate}

\end{document}